\newpage\SourceLeaf{57}\ArithmeticHead{50}{L}{BACHMANN, ÜBER GAUSS’ ZAHLENTHEORETISCHE ARBEITEN.}
\begin{center}\bfseries 20.\end{center}
\noindent\hspace*{\parindent}Werfen wir noch einen Blick auf die Chronologie der acht \textsc{Gauss}schen Beweise für das Fundamentaltheorem. Datieren wir sie nach der Zeit ihrer Veröffentlichung, und zählen sie, wie es \textsc{Gauss} in seinen Publikationen tut, dementsprechend, so ist
\begin{center}\small\setlength{\tabcolsep}{3pt}
\begin{tabular}{@{}rllccr@{}}
Beweis & I & (D. A. art. 135 sqq, W. I, S. 104) & vom & Jahre & 1801\\
„ & II & (ebendas. art. 262, W. I, S. 292) & „ & „ & 1801\\
„ & III & (Comm. Gotting. \textbf{16}, W. II, S. 1) & „ & „ & 1808\\
„ & IV & (Comm. Gott. rec. \textbf{1}, W. II, S. 9) & „ & „ & 1811\\
„ & V u. VI & (Comm. Gott. rec. \textbf{4}, W. II, S. 47) & „ & „ & 1818\\
„ & VII u. VIII & (Nachlass, W. II, S. 234) & „ & „ & 1863.
\end{tabular}\end{center}
\noindent Von diesen Beweisen ist I der zeitlich \textls[100]{erste} und am 8. April 1796 (Nr. 2 T.) von \textsc{Gauss} gefunden; Beweis II ist der \textls[100]{zweite}, am 27. Juni 1796 (Nr. 16 T.) gefunden, die Beweise VII und VIII, die im Grunde nur einen vollen Beweis bilden, sind spätestens am 2. September 1796 (Nr. 30 T.) von \textsc{Gauss} erhalten und in seinem Nachlass als \textls[100]{dritter und vierter} (W. II, S. 234) ausdrücklich bezeichnet. Der Beweis IV ist schon Mitte Mai 1801 (Nr. 118 T.) von \textsc{Gauss} gefunden und dort als \textls[100]{fünfter} gezählt worden. Die übrigen Beweise entstammen, wie bemerkt, nach seiner eigenen Aussage seinen Bemühungen um die kubischen und biquadratischen Reste, die erst 1805 einsetzten; sicherlich ist die in Nr. 134 T. auf den 6. Mai 1807 datierte »demonstratio principiis omnino elementaribus innixa« der Beweis III. Ein Anfang desselben, insbesondere die oben angegebene Formel für den Dezidenten \((q,p)\) findet sich schon im Handbuche (18, Bd) vom Oktober 1805 des \textsc{Gauss}schen Nachlasses S. 164 (abgedruckt W. X1, S. 26) vor der Notiz über den Beweis der Irreduzibilität der allgemeinen Kreisteilungsgleichung, also vor dem 12. Juni 1808.

In einem andern im September 1813 begonnenen Handbuche (21, Bg) steht S. 6—8 ein Aufsatz mit dem Titel »\textit{Dritter Beweis des Fundamentaltheorems bei den quadratischen Resten in einer neuen Einkleidung}« und mit dem Datum Novb. 12, jedenfalls 1813 (abgedruckt W. X1, S. 33), der mit den Worten endigt: »Bei dieser Einkleidung des Beweises ist der wahre Nerf desselben mehr in die Augen fallend als bei derjenigen, in welcher er in den Göttingischen Commentationen Bd. XVI erscheint«. Da dieser Aufsatz im Prinzip
\newpage\SourceLeaf{58}\ArithmeticHead{51}{R}{DIE SPÄTEREN BEWEISE DES FUNDAMENTALTHEOREMS.}\HistoricalFoot{}{7*}
\noindent mit dem Beweise V identisch und, nur in der Entwicklung desselben von ihm abweichend, eine blosse Modifikation desselben ist, so ist ersichtlich auch Beweis V erst durch eine vereinfachende Neubearbeitung aus Beweis III hervorgegangen, also später als dieser. Da aber (W. II, S. 50) die Beweise V und VI als die vor neun Jahren versprochenen, nämlich (W. II, S. 43) als die bei Veröffentlichung von Beweis IV (24. August 1808) schon vorhandenen Beweise bezeichnet werden, so muss Beweis V zwischen dem 6. Mai 1807 und dem 24. August 1808 entstanden sein. Im \textit{Tagebuch} geschieht seiner keine Erwähnung, ebensowenig des Beweises VI, dessen Datierung am unsichersten bleibt.

Beim Beweise III (W. II, S. 4) und in der Anzeige desselben (W. II, S. 153) erwähnt \textsc{Gauss} drei Beweise, die er nach dem Beweise I gefunden und welche »sehr tiefliegende und ihrem Inhalte nach ganz heterogene Untersuchungen voraussetzen«; einer davon sei der Beweis II. Von den beiden übrigen ist unter dem einen zweifelsohne der Beweis IV gemeint; der andere kann Beweis V nicht sein, der, wie gesagt, später als Beweis III ist, und auf welchen die erwähnte Charakterisierung nicht passt. Es bleibt also nur die Wahl zwischen dem Doppelbeweise der A. R. (der demonstratio tertia und quarta, welche eigentlich erst zusammen eine einzige »demonstratio completa« ausmachen) und dem — ihm übrigens nahestehenden — Beweise VI, und man möchte sich \textsc{Kronecker} entgegen für den erstern entscheiden, da es wunderbar scheinen muss, dass \textsc{Gauss} diesen mit Stillschweigen übergangen haben sollte. Jenachdem man nun sein Zitat auf den Beweis VI bezieht oder nicht, würde dieser Beweis vor Beweis III d. i. vor den 6. Mai 1807 oder, wie Beweis V, zwischen diesen Zeitpunkt und den 24. August 1808 zu datieren sein; ob dann vor den Beweis V oder nach denselben, muss dahingestellt bleiben.

Der deutschen Darstellung des Beweises V geht in dem Handbuche (21, Bg) vom September 1813 auf S. 4, 5 auch eine (im wesentlichen mit der lateinischen übereinstimmende) deutsche Darstellung des Beweises VI vorauf (abgedruckt W. X1, S. 28), die als \textls[100]{fünfter} Beweis gezählt wird, umgekehrt wie in den Commentationen, wo er dem Beweise V folgt und als Demonstratio sexta bezeichnet wird. Da die lateinischen Fassungen dieser Beweise in den Comm. Gott. rec. \textbf{4} von \textsc{Gauss} erst 1817 der Öffentlichkeit übergeben
\newpage\SourceLeaf{59}\ArithmeticHead{52}{L}{BACHMANN, ÜBER GAUSS’ ZAHLENTHEORETISCHE ARBEITEN.}
\noindent sind, möchte man annehmen, dass die deutschen Darstellungen die ursprüngliche Fassung der schon 1808 vorhanden gewesenen Beweise wiedergeben. Der Algorithmus zur Bestimmung des quadratischen Charakters einer Zahl \(B\) in bezug auf eine andere \(A\), mit welchem die dritte Abhandlung schliesst, kann hinwiederum nicht wohl vor dem August 1808 gefunden sein, da er in dem Handbuche vom Oktober 1805 erst S. 213 nach astronomischen Beobachtungen von jenem Datum aufgeführt wird.
